\documentclass[11pt,a4paper]{article}
\usepackage[T1]{fontenc}\usepackage[utf8]{inputenc}\usepackage{lmodern}\usepackage{microtype}
\usepackage[a4paper,margin=2.25cm]{geometry}\usepackage{longtable,booktabs,array,calc}\usepackage{fancyvrb}
\usepackage{xcolor}\usepackage{hyperref}\usepackage{xurl}\usepackage{fancyhdr}\usepackage{enumitem}\usepackage{parskip}\usepackage{titlesec}
\definecolor{ddtadark}{RGB}{30,38,48}\definecolor{ddtablue}{RGB}{38,82,126}
\hypersetup{colorlinks=true,linkcolor=ddtablue,urlcolor=ddtablue,citecolor=ddtablue}
\pagestyle{fancy}\fancyhf{}\lhead{Documentation-Driven Threat Analysis}\rhead{BA3 identity/lifecycle candidate}\cfoot{\thepage}\setlength{\headheight}{14pt}
\setlist[itemize]{leftmargin=1.5em,itemsep=0.2em,topsep=0.25em}\setlist[enumerate]{leftmargin=1.7em,itemsep=0.2em,topsep=0.25em}
\titleformat{\section}{\Large\bfseries\color{ddtadark}}{}{0em}{}\titleformat{\subsection}{\large\bfseries\color{ddtadark}}{}{0em}{}\titleformat{\subsubsection}{\normalsize\bfseries\color{ddtadark}}{}{0em}{}
\providecommand{\tightlist}{\setlength{\itemsep}{0pt}\setlength{\parskip}{0pt}}\setlength{\emergencystretch}{4em}\hbadness=10000\hfuzz=20pt\sloppy
\begin{document}
\section{DDTA BA3 cross-baseline identity/lifecycle candidate -
R1}\label{ddta-ba3-cross-baseline-identitylifecycle-candidate---r1}

\textbf{Status:} PROVISIONAL CANDIDATE AFTER BA3-T2 / NOT CLOSED / BA3
OPEN

\textbf{Derived by:} BA3-T2 cross-baseline identity, staleness and
lifecycle pressure test

\textbf{Repository baseline reviewed:}
\texttt{52864d2ce177abdd694436306d8152db688effa0}

\textbf{Closed dependencies:} BA0 responsibility boundary; BA1
\texttt{BAReferent\ +\ BAProposition}; BA2 relation/action vocabulary;
BA3-T1 provenance/origin lower bound.

\subsection{1. Purpose}\label{purpose}

BA3-T2 asks what must remain stable, what must be revalidated and what
must be replaced when governed documentation changes between baselines.

The pressure is mutation-driven. The candidate is not a generic
version-control model and does not make Git history, document revision
events or ThreatForge runtime state part of Base Analysis semantics.

The core result is that \textbf{baseline provenance, semantic identity
continuity, review state and cross-baseline disposition are distinct
concerns}.

\subsection{2. Candidate cross-baseline lower
bound}\label{candidate-cross-baseline-lower-bound}

BA3-T1 already requires one provenance/origin attachment per
materialized BA identity and baseline. BA3-T2 adds two
representation-independent metadata contracts:

\begin{verbatim}
BABaselineReview
|- targetElement          1     BAReferent | BAProposition
|- evaluatedBaselineKey   1
|- reviewState            1     PENDING_REVIEW | ACCEPTED | REJECTED
`- freshness              1     CURRENT | STALE

BACrossBaselineContinuity
|- priorElement           1     BAElementRef@priorBaseline
|- targetBaselineKey      1
|- disposition            1     RETAIN | REPLACE | RETIRE
|- successorElement       0..1  required for REPLACE
`- continuityBasis        1..*  GovernedSourceRef | BAElementRef
\end{verbatim}

These are metadata contracts, not BAE identity families and not
mandatory physical records/tables.

\texttt{SUPERSEDED} and \texttt{RETIRED} are lifecycle interpretations
of an accepted continuity disposition:

\begin{verbatim}
accepted REPLACE -> prior element is SUPERSEDED
accepted RETIRE  -> prior element is RETIRED
\end{verbatim}

A retained element keeps the same semantic identity across the reviewed
baselines. A newly introduced element has no predecessor continuity
obligation.

\subsection{3. Why the dimensions must remain
separate}\label{why-the-dimensions-must-remain-separate}

BA3-T1 already closed:

\begin{verbatim}
originState = GROUNDED | DERIVED | DIAGNOSTIC_UNRESOLVED
\end{verbatim}

That dimension answers \textbf{where/how meaning originated}.

BA3-T2 requires a separate review dimension because all of the following
are valid:

\begin{verbatim}
GROUNDED              + ACCEPTED
DERIVED               + ACCEPTED
DERIVED               + REJECTED
DIAGNOSTIC_UNRESOLVED + ACCEPTED
  [accepted as a diagnostic, not as project truth]
GROUNDED              + STALE       [relative to a newer baseline pending review]
\end{verbatim}

Therefore:

\begin{verbatim}
originState != reviewState != freshness/continuity disposition
\end{verbatim}

Collapsing them would either treat a diagnostic as rejected merely
because it is unresolved, or treat a grounded statement as current
merely because it once came directly from a governed source.

\subsection{\texorpdfstring{4. \texttt{STALE} is target-baseline
uncertainty, not historical
invalidity}{4. STALE is target-baseline uncertainty, not historical invalidity}}\label{stale-is-target-baseline-uncertainty-not-historical-invalidity}

An accepted BA element from baseline \texttt{B0} does not become
historically false merely because \texttt{B1} exists.

\texttt{STALE} means:

\begin{quote}
the prior accepted meaning is potentially affected in the target
baseline and MUST NOT be treated as current for that target baseline
until revalidated.
\end{quote}

Thus:

\begin{verbatim}
valid for B0
    can coexist with
STALE relative to B1 review
\end{verbatim}

After review, the target-baseline outcome is one of:

\begin{itemize}
\tightlist
\item
  \texttt{RETAIN} -\textgreater{} same identity is accepted/current in
  B1;
\item
  \texttt{REPLACE} -\textgreater{} a successor identity is
  accepted/current and the old element is superseded;
\item
  \texttt{RETIRE} -\textgreater{} no successor is required and the old
  element is retired;
\item
  rejection of a proposed successor -\textgreater{} it does not enter
  the accepted B1 Base Analysis.
\end{itemize}

\texttt{STALE} is intentionally not a project-document lifecycle state.
It is an analysis-layer revalidation state.

\subsection{\texorpdfstring{5. Cross-baseline identity rule for
\texttt{BAReferent}}{5. Cross-baseline identity rule for BAReferent}}\label{cross-baseline-identity-rule-for-bareferent}

A \texttt{BAReferent} identity is retained when the same independently
reusable methodology-neutral project meaning survives the governed
change.

Changes in the following do \textbf{not by themselves} force a new
referent identity:

\begin{itemize}
\tightlist
\item
  source wording;
\item
  source locator;
\item
  realization technology;
\item
  ownership/responsibility relations;
\item
  provider placement;
\item
  propositions that classify, constrain or relate the referent.
\end{itemize}

A referent is replaced or retired when the independently reusable
meaning itself changes or ceases to exist.

Candidate test:

\begin{quote}
If consumers can still point to the same project-semantic thing across
the two baselines while truthfully changing propositions about it,
retain the referent identity. If the supposed continuity would make two
materially different project meanings appear to be one thing, replace
it.
\end{quote}

\subsection{\texorpdfstring{6. Cross-baseline identity rule for
\texttt{BAProposition}}{6. Cross-baseline identity rule for BAProposition}}\label{cross-baseline-identity-rule-for-baproposition}

\texttt{BAProposition} has a stricter continuity test because its
identity is assertion identity.

A proposition may be retained only when its normalized assertion meaning
remains equivalent across baselines. For the closed BA2 structure this
requires materially equivalent:

\begin{verbatim}
semanticOperatorKey
polarity
role-bound participants / controlled local values
condition / temporalScope modifiers
\end{verbatim}

Participants are compared through accepted referent continuity, not
lexical labels.

Changes only to provenance locators, source wording or display labels do
not force proposition replacement if the assertion meaning is unchanged.

If a meaning-bearing operator, polarity, participant binding, governed
constraint value or local modifier changes, the old proposition MUST NOT
be silently edited under the same assertion identity. It is
replaced/superseded or retired.

This prevents \texttt{BAProposition} from becoming a mutable document
slot.

\subsection{7. Source-document revision does not determine BA
continuity}\label{source-document-revision-does-not-determine-ba-continuity}

A governed document can be revised while some BA elements are retained
and others are replaced.

Therefore the source-layer dispositions:

\begin{verbatim}
retain / revise / replace / retire
\end{verbatim}

MUST NOT be copied mechanically into BA identity disposition.

In particular, BA3-T2 does \textbf{not} require \texttt{REVISE} as a
separate BA semantic-identity disposition. A source revision yields one
of three BA outcomes after semantic review:

\begin{verbatim}
same meaning survives      -> RETAIN
assertion/meaning replaced -> REPLACE
meaning disappears         -> RETIRE
\end{verbatim}

Non-identity metadata/provenance may of course be refreshed on a
retained identity.

\subsection{8. Mutation M1 - Ethernet to
Wi-Fi}\label{mutation-m1---ethernet-to-wi-fi}

The governed mutation changes only \texttt{D-3.5} from wired Ethernet to
available Wi-Fi while transport remains externally consumed and
camera/processor remain separated.

\subsubsection{8.1 Referent continuity}\label{referent-continuity}

\begin{verbatim}
LocalConnectivity       RETAIN
CameraSubsystem         RETAIN
RecognitionProcessor    RETAIN
RecognitionCapture      RETAIN
WiredEthernet           RETIRE if no other governed use remains
WiFi                    INTRODUCE
\end{verbatim}

The abstract connectivity meaning survives even though its realization
changes.

\subsubsection{8.2 Proposition continuity}\label{proposition-continuity}

\begin{verbatim}
realize(LocalConnectivity, WiredEthernet)
    REPLACE
        -> realize(LocalConnectivity, WiFi)

transfer(CameraSubsystem,
         RecognitionProcessor,
         RecognitionCapture)
    RETAIN after review

assignResponsibility[negative](project,
                               transport,
                               ownership/management)
    RETAIN after review
\end{verbatim}

The old realization proposition is therefore \texttt{SUPERSEDED}, not
merely textually revised. The source \texttt{D-3.5} may be revised under
the same document identity while BA assertion identity changes.

The transfer and security constraints remain current only \textbf{after
review}; their survival is not inferred solely from unchanged files.

\subsection{9. Mutation M2 - external transport to project-owned
transport}\label{mutation-m2---external-transport-to-project-owned-transport}

\texttt{D-3.4} changes the responsibility boundary from consuming
external transport to project ownership/management.

\subsubsection{9.1 Meanings that can
survive}\label{meanings-that-can-survive}

\begin{verbatim}
LocalConnectivity / LocalTransportCapability
  RETAIN if the abstract capability meaning is preserved
RecognitionCapture                              RETAIN
transfer delivery behavior                      RETAIN
confidentiality/integrity/provenance constraints RETAIN after review
\end{verbatim}

\subsubsection{9.2 Responsibility
propositions}\label{responsibility-propositions}

The negative responsibility assertion cannot remain the same proposition
when polarity changes:

\begin{verbatim}
assignResponsibility [negative]
  project -> transport -> ownership/management

    REPLACE

assignResponsibility [affirmative]
  project -> transport -> ownership/management
\end{verbatim}

The prior proposition is \texttt{SUPERSEDED}.

The old external \texttt{consumeService} relation is normally
\texttt{RETIRED} if the project no longer consumes that transport as an
external service. A new project-owned transport FR branch may introduce
new referents and propositions; this does not require rewriting the
consumer's existing transfer/security semantics.

\subsection{10. Mutation M3 - remote recognition to local
recognition}\label{mutation-m3---remote-recognition-to-local-recognition}

\texttt{D-3.3} changes structurally from separate recognition processor
to recognition on the camera.

\subsubsection{10.1 Referent outcomes}\label{referent-outcomes}

\begin{verbatim}
CameraSubsystem          RETAIN
RecognitionCapture       RETAIN if acquisition/evidence meaning remains
RecognitionProcessor     REVIEW; RETIRE if it has no remaining governed role
remote delivery behavior RETIRE
\end{verbatim}

The same capture identity can survive even when the transfer relation
disappears. This is direct evidence that referent and proposition
lifecycle rules cannot be identical.

\subsubsection{10.2 Proposition outcomes}\label{proposition-outcomes}

The remote transfer assertion no longer has a true successor with the
same project meaning:

\begin{verbatim}
transfer(CameraSubsystem,
         RecognitionProcessor,
         RecognitionCapture)
    RETIRE
\end{verbatim}

It is not rewritten into an unrelated local-processing assertion.

Constraints whose target is specifically the retired delivery behavior
are also retired in their old form unless governed documentation
explicitly establishes successor obligations. They MUST NOT be dragged
automatically onto local recognition merely because the concern remains
security-relevant.

A placement/responsibility proposition for \texttt{D-3.3} is replaced by
the new local-placement proposition if a semantically corresponding
successor is materialized.

\texttt{D-3.4}/\texttt{D-3.5}-derived BA elements become at least
\textbf{STALE for review} where their relevance depended on the removed
capture transfer. They are not automatically retired because other
project uses may remain.

\subsection{11. Order-fulfillment responsibility-boundary
control}\label{order-fulfillment-responsibility-boundary-control}

The order corpus chooses project-owned internal inventory authority and
explicitly states that an external WMS would move many project-owned
receipt/adjustment/availability/reservation/concurrency/release/issue
obligations toward integration/contract semantics.

The counterfactual applies the same rules:

\begin{itemize}
\tightlist
\item
  stable project-semantic meanings such as \texttt{OrderEvaluation},
  governed \texttt{ReservationResult} and externally visible reservation
  semantics may retain identity where the contract preserves them;
\item
  internal \texttt{InventoryService}/\texttt{InventoryLedger} meanings
  and internal mutation propositions are retired or replaced if project
  responsibility disappears;
\item
  external WMS/service/contract meanings and new
  \texttt{consumeService}/normalization propositions are introduced;
\item
  source-level responsibility change does not force wholesale
  replacement of unrelated order/payment/fulfillment BA identities.
\end{itemize}

This cross-corpus control supports the family-specific continuity rules
without introducing domain-specific lifecycle types.

\subsection{12. Review state is required separately from origin
state}\label{review-state-is-required-separately-from-origin-state}

BA3-T2 finds that \texttt{PENDING\_REVIEW}, \texttt{ACCEPTED} and
\texttt{REJECTED} are required analytical review states at the
element/materialization boundary.

Why:

\begin{enumerate}
\def\labelenumi{\arabic{enumi}.}
\tightlist
\item
  a prior accepted element may become stale relative to a new baseline
  before review;
\item
  a tool/analyst-derived candidate must not silently become accepted BA
  meaning;
\item
  an accepted diagnostic must remain distinguishable from a rejected
  derivation;
\item
  replacement candidates may be rejected without rewriting the
  historical predecessor.
\end{enumerate}

A rejected proposal may be retained for audit/history, but it is not
part of the accepted Base Analysis for that baseline.

This state is analysis-layer governance. It is not governed-document
acceptance and does not make BA project authority.

\subsection{13. Diagnostic resolution across
baselines}\label{diagnostic-resolution-across-baselines}

A diagnostic with:

\begin{verbatim}
originState = DIAGNOSTIC_UNRESOLVED
\end{verbatim}

MUST NOT be retyped in place as \texttt{GROUNDED} merely because the
source was later corrected.

If the governed correction removes the ambiguity/conflict/missing basis:

\begin{verbatim}
old diagnostic -> RETIRE
\end{verbatim}

If a materially different unresolved problem remains:

\begin{verbatim}
old diagnostic -> REPLACE -> new diagnostic
\end{verbatim}

Any new or retained grounded/derived project-semantic elements in the
corrected baseline have their own origin/provenance and review state.

This preserves historical truth: the earlier diagnostic really existed
and was unresolved in its baseline.

\subsection{14. Minimal staleness
triggers}\label{minimal-staleness-triggers}

BA3-T2 does not close a full change-impact algorithm, but mutation
pressure forces a minimum rule for deciding what requires review in a
target baseline.

A previously accepted BA element is at least a \textbf{staleness
candidate} when:

\begin{enumerate}
\def\labelenumi{\arabic{enumi}.}
\tightlist
\item
  one of its direct \texttt{sourceLink} targets changes or disappears;
\item
  one of its \texttt{derivationBasis} elements/sources becomes stale,
  replaced or retired;
\item
  one of a proposition's role-bound referents is replaced or retired;
\item
  a governed dependency/effective context known to justify the element
  changes materially.
\end{enumerate}

Staleness may propagate transitively through derivation basis. It MUST
NOT propagate to the entire Base Analysis merely because one document
file changed.

The exact dependency closure, effective-governed-context representation
and re-analysis scheduling remain later work.

\subsection{15. Accepted provisional lifecycle
semantics}\label{accepted-provisional-lifecycle-semantics}

\begin{verbatim}
PENDING_REVIEW
  candidate or impacted carry-forward not yet accepted for target baseline

ACCEPTED
  accepted analytical materialization for the evaluated baseline

REJECTED
  reviewed candidate not admitted to accepted Base Analysis

STALE
  prior accepted meaning potentially impacted in target baseline
  and not yet revalidated

SUPERSEDED
  prior accepted element has an accepted REPLACE successor

RETIRED
  prior accepted element no longer applies and no successor identity is required
\end{verbatim}

\texttt{SUPERSEDED} and \texttt{RETIRED} preserve history; neither
deletes the old origin/provenance record.

\subsection{16. Falsification rules}\label{falsification-rules}

Revise this candidate if a concrete governed corpus demonstrates that:

\begin{enumerate}
\def\labelenumi{\arabic{enumi}.}
\tightlist
\item
  \texttt{BAReferent} and \texttt{BAProposition} can use one identical
  identity-continuity rule without semantic distortion;
\item
  a proposition can change operator/polarity/role-bound meaning while
  honestly preserving assertion identity;
\item
  source-layer \texttt{revise} maps one-to-one to BA identity revision
  without ambiguity;
\item
  staleness can be inferred from origin state alone;
\item
  accepted/rejected review state can be removed without permitting
  silent candidate promotion;
\item
  a resolved diagnostic can safely mutate its historical origin from
  unresolved to grounded; or
\item
  a responsibility-boundary mutation forces a new BA1/BA2 semantic
  family rather than lifecycle/provenance handling.
\end{enumerate}

No reviewed evidence currently forces such a revision.

\subsection{17. Provisional
dispositions}\label{provisional-dispositions}

\begin{verbatim}
Stable identity across baselines                       REQUIRED CAPABILITY
BAReferent continuity rule
  RETAIN WHEN REFERENT MEANING SURVIVES
BAProposition continuity rule
  STRICT SEMANTIC EQUIVALENCE REQUIRED
Source document revision == BA identity revision       REJECTED
BA continuity disposition                              RETAIN | REPLACE | RETIRE
BA-level REVISE as fourth identity disposition         NOT REQUIRED
PENDING_REVIEW / ACCEPTED / REJECTED                   REQUIRED
CURRENT / STALE distinction                            REQUIRED
REPLACE -> SUPERSEDED
  REQUIRED SEMANTIC INTERPRETATION
RETIRE -> RETIRED
  REQUIRED SEMANTIC INTERPRETATION
Resolved diagnostic mutates to GROUNDED                REJECTED
Historical provenance deletion on replacement          REJECTED
Third BAE lifecycle family                              NOT FORCED
BA1 reopen                                              NOT TRIGGERED
BA2 reopen                                              NOT TRIGGERED
BA3                                                     STARTED / NOT CLOSED
\end{verbatim}

\subsection{18. Remaining BA3 questions}\label{remaining-ba3-questions}

The smallest unresolved set after T2 is:

\begin{enumerate}
\def\labelenumi{\arabic{enumi}.}
\tightlist
\item
  exact method-neutral derivation-rule/rationale registry and
  reproducibility contract;
\item
  effective governed context and dependency representation sufficient
  for principled staleness/change-impact propagation;
\item
  source-to-analysis and analysis-to-source feedback lineage after a
  corrective governed change;
\item
  closure review over provenance + identity/lifecycle once those
  dependencies are pressure-tested.
\end{enumerate}

The next microstep should attack derivation reproducibility and
change-impact lineage without designing BA4 or analysis-method schemas.

\end{document}
